\documentclass[10pt,a4paper]{amsart}
\usepackage[margin=19mm]{geometry}
\usepackage{amsmath,amssymb,mathtools}
\usepackage[scr=rsfs,cal=euler]{mathalfa}
\usepackage{xfrac}
\usepackage[unicode,hidelinks,bookmarks=false]{hyperref}
\newcommand{\ie}{i.e.\ }
\newcommand{\cf}{cf.\ }
\newcommand{\resp}{respectively\ }
\newcommand{\Subsection}{\subsection}
\providecommand{\nobreakdash}{\nobreak\hbox{-}}
\setlength{\emergencystretch}{2em}
\multlinegap0pt
\usepackage{bm}
\usepackage{bbm}
\usepackage{dsfont}
\usepackage{xspace}
\usepackage{cancel}
\usepackage{mathtools}
\usepackage{amsthm}
\makeatletter
\@ifundefined{thm}{\newtheorem{thm}{Thm}}{}
\@ifundefined{lem}{\newtheorem{lem}[thm]{Lemma}}{}
\@ifundefined{prop}{\newtheorem{prop}[thm]{Proposition}}{}
\makeatother
\def\bbF{\mathbb F} 
\def\bbK{\mathbb K} 
\newcommand{\Hom}{\operatorname{Hom}}
\newcommand{\Ker}{\operatorname{Ker}}
\newcommand{\Rep}{\operatorname{Rep}}
\newcommand{\lto}{\longrightarrow}
\newcommand{\rep}{\operatorname{rep}}
\newcommand{\xto}{\xrightarrow}
\title[Multisets of finite intervals and a universal category of po]{Multisets of finite intervals and a universal category of poset representations}
\author{Henning Krause, Balduin Stoye}
\date{Source: arXiv:2601.22649v2 (CC BY 4.0)}
\begin{document}
\maketitle

\noindent\emph{Source and licence.} Multisets of finite intervals and a universal category of poset representations. Version arXiv:2601.22649v2 (CC BY 4.0), distributed under the Creative Commons Attribution 4.0 International licence, as recorded in the arXiv licence metadata for this exact version and shown on the abstract page https://arxiv.org/abs/2601.22649v2. Full rights notice: licenses/arxiv-2601.22649-CC-BY-4.0.txt.\par\smallskip

\noindent\textit{Editorial scope.} Counted unit: the Prop whose source label is \texttt{pr:coherent} in the article (source0.tex lines 1162--1171, source label \texttt{pr:coherent}), delivered together with the article's own complete original proof of it (source0.tex lines 1173--1218, one \texttt{proof} environment of 46 lines). No mathematics is written, repaired or re-derived: statement and proof are the article's own text line for line. Required context retained uncounted: eq:tensor (lines 1027--1030); le:mod-coherent (lines 1141--1148). Every definition, display and label of those lines is retained unchanged. Omitted from this package and not redistributed: 1--1026, 1031--1140, 1149--1161, 1172--1172, 1219--1324. Nothing inside a retained excerpt is omitted or altered: every retained body is delivered exactly as the article's lines are written, so no omission receipt of any kind accompanies this package. Citations: the retained excerpts carry no citation command at all, so no bibliography entry of the article is transcribed and no reference is redistributed. Reference adaptation: none was needed and none was made; every reference inside the delivered unit resolves inside it, so no unresolved reference or citation is produced. Printed slips kept literal: no printed word, symbol, hypothesis or number of the counted statement or of its proof was altered. Layout and class: the article's own class is \texttt{amsart} with packages including amsmath, graphicx, amssymb, mathrsfs, nicematrix, stmaryrd, euscript, tikz, tikz-cd, hyperref; the wrapper is the lane's pinned \texttt{amsart} editorial wrapper at 10pt on A4, which restates the theorem-like environments the retained text needs and adds the article's own macro definitions that the retained text needs (\texttt{\textbackslash Hom}, \texttt{\textbackslash Ker}, \texttt{\textbackslash Rep}, \texttt{\textbackslash bbF}, \texttt{\textbackslash bbK}, \texttt{\textbackslash lto}, \texttt{\textbackslash rep}, \texttt{\textbackslash xto}), with extra wrapper packages (bm, bbm, dsfont, xspace, cancel, mathtools). The delivered numbering is the unit's own. Assets: no retained span contains a figure environment, a graphics-inclusion command, a PDF-page inclusion, a table or a TikZ picture; no image file is copied into this package, the package declares no asset dependency and no image file is redistributed with it. Licence: version arXiv:2601.22649v2 of this article is distributed under the Creative Commons Attribution 4.0 International licence, as recorded in the arXiv licence metadata for this exact version and shown on the abstract page https://arxiv.org/abs/2601.22649v2. A full rights notice is delivered at licenses/arxiv-2601.22649-CC-BY-4.0.txt. Characters: every retained excerpt is pure ASCII, so the delivered engine, XeLaTeX with its default Latin Modern text font, reports no missing character.
\begin{equation}\label{eq:tensor}
  \Hom_{\bbK P}(F\otimes_\bbK M,G)\cong\Hom_{\bbK}(M,\Hom_{\bbK
    P}(F,G)).
\end{equation}

\begin{lem}\label{le:mod-coherent}
 For a Grothendieck category with a generating set of objects that are
  finitely generated and projective, the following  are equivalent.
  \begin{enumerate}
\item All finitely generated projective objects are coherent.
\item The finitely presented objects form an abelian category.
 \end{enumerate}
\end{lem}

\begin{prop}\label{pr:coherent}
  For a poset $P$ the following are equivalent.
\begin{enumerate}
\item The poset $P$ is coherent.
\item For each $x\in P$ the compact elements in the lattice of ideals
  of $P$ contained in $\downarrow\! x$ are closed under finite meets.
\item For each $x\in P$ the object $h_x\in\Rep(P,\bbF)$ is coherent.
\item The category $\rep(P,\bbF)$ is abelian.
\end{enumerate}
\end{prop}

\begin{proof}
  (1) $\Leftrightarrow$ (2): The ideals of a poset form a distributive
  lattice, and an ideal $I$ is compact if and only if
  $I=\bigcup_\alpha\downarrow\! x_\alpha$ for some finite set of elements $x_\alpha$.
  Thus to check condition (2) it suffices to show that $\downarrow\!
  y\cap {\downarrow\! y'}$ is compact for each pair $y,y'\le x$, and
  this is precisely condition (1).

  (2) $\Rightarrow$ (3): For a morphism
  $\phi\colon \bigoplus_{\alpha\in A} h_{y_\alpha}\to h_x$ given by a finite set
  of elements $y_\alpha\in {\downarrow\! x}$ we need to show that its
  kernel is finitely generated.  For $z\in P$ set
  \[V_z:=\Ker\left( \bigoplus_{\alpha\in A} h_{y_\alpha}(z)\to
      h_x(z)\right).\] The inclusion
  \[V_z\lto \bigoplus_{\alpha\in A} h_{y_\alpha}(z)=\Hom_{\bbF P}(h_z,
    \bigoplus_{\alpha\in A} h_{y_\alpha})\] corresponds via \eqref{eq:tensor}
  to a morphism
  \[\psi_z\colon h_z\otimes_\bbF V_z\lto \bigoplus_{\alpha\in A}
  h_{y_\alpha}.\] Using the adjunction \eqref{eq:tensor} it follows that
  $\phi\psi_z=0$ and any morphism
  $\psi\colon h_z\to \bigoplus_{\alpha\in A} h_{y_\alpha}$ satisfying
  $\phi\psi=0$ factors through $\psi_z$.

  Using the coherence of $P$ there is for each subset $B\subseteq A$ a
  finite subset $B'\subseteq P$ satisfying
\[\bigcap_{\alpha\in B}\downarrow\!
  y_\alpha=\bigcup_{z\in B'}\downarrow\! z.\]
This yields a morphism
\[\bigoplus_{B\subseteq A}\bigoplus_{z\in B'}
  h_z\otimes_\bbF V_z\xto{\,(\psi_z)\,} \bigoplus_{\alpha\in A}
  h_{y_\alpha},\] where $B\subseteq A$ runs through all subsets, and its
image equals $\Ker\phi$ by construction. In fact, for any
$\psi\colon h_z\to\Ker\phi$ choose $B=\{\alpha\in A\mid\psi_\alpha\neq 0\}$. Then
$z\le z'$ for some $z'\in B'$ and $\psi$ factors through $\psi_{z'}$.

(3) $\Rightarrow$ (1): A diagram $y\le x\ge
y'$ yields a morphism $h_y\oplus h_{y'}\to h_x$, and its kernel is
given by diagrams $y\ge z\le y'$. Suppose the kernel is
a  quotient of $\bigoplus_\alpha h_{z_\alpha}$  for some finite set of elements
  $z_\alpha\in\downarrow\! x$.  This yields finitely many diagrams $y\ge z_\alpha\le y'$ such that for each
diagram $y\ge z\le y'$ there is an index $\alpha$ with  $z\le z_\alpha$. 

(3) $\Leftrightarrow$ (4): This follows from
Lemma~\ref{le:mod-coherent}, once we observe that coherent objects are
closed under finite direct sums and summands.
\end{proof}

\end{document}
